%macbook 
\loai{Số vân sáng tại M khi biết rõ toạ độ}
\begin{vidu}%[3P5-6.3K][Loai1] 
Trong thí nghiệm giao thoa ánh sáng của Y-âng có $a=1$ mm, $D=1\ m$, ánh sáng thí nghiệm là ánh sáng trắng có bước sóng từ $0,4\ \mu m$ đến $0,75\ \mu m$. Tại điểm M cách vân trung tâm 5 mm có mấy quang phổ chồng lên nhau?
\choice
{5}
{4}
{\True 6}
{7}
\loigiai{
\begin{itemize}
\item Giả sử tại M có vân sáng của bức xạ $\lambda $: $x=k.\dfrac{\lambda D}{a}\Rightarrow \lambda =\dfrac{a.x}{kD}=\dfrac{1.10^{-3}{.5.10}^{-3}}{k.1}=\dfrac{5}{k}\ \ \mu m $.
\item Ánh sáng thí nghiệm là ánh sáng trắng có bước sóng từ $0,4\ \mu m$ đến $0,75\ \mu m$ nên:
\begin{align*}
0,4\le \dfrac{5}{k}\le 0,75\Rightarrow 6,7\le k\le 12,5\Rightarrow k=\left\{{7;\ 8;\ 9;\ 10;\ 11;\ 12}\right\}.
\end{align*}
\item Cứ một giá trị k, ứng với nó là một bức xạ cho vân sáng tại M.
\item Vậy, tại M có tổng cộng 6 vân sáng của 6 bức xạ chồng lên nhau.
\end{itemize}
}
\end{vidu} \haidong{20}
\loai{Số vân sáng tại M khi phải tìm toạ độ}
\begin{vidu}%[3P5-6.3K][Loai2]
Trong thí nghiệm giao thoa Y-âng thực hiện với ánh sáng trắng có bước sóng từ $0,38\ \mu m$ đến $0,76\ \mu m$. Có bao nhiêu bức xạ đơn sắc cho vân sáng trùng vân sáng bậc 3 của bức xạ có bước sóng $0,76\ \mu m$?
\choice
{2}
{\True 3}
{4}
{5}
\loigiai{ 
\begin{itemize}
\item Ta có: $x_M=k\dfrac{\lambda D}{a}\Rightarrow \lambda =\dfrac{ax_M}{kD}=\dfrac{2,28}{k}\ \mu m\xrightarrow{{0,38\leq \lambda =\tfrac{2,28}{k}\leq 0,76}}3\leq k\leq 6\Rightarrow k=4;5;6$.
\item 
\item 
\end{itemize}
}
\end{vidu} \haidong{20}

\loai{Tìm chính xác bức xạ cho vân sáng tại M}
\begin{vidu}%[3P5-6.3G][Loai3]
\nguon{ĐH - 2010} Trong thí nghiệm Y-âng về giao thoa ánh sáng, hai khe được chiếu bằng ánh sáng trắng có bước sóng từ 380 nm đến 760 nm. Khoảng cách giữa hai khe là 0,8\ mm, khoảng cách từ mặt phẳng chứa hai khe đến màn quan sát là 2 m. Trên màn, tại vị trí cách vân trung tâm 3\ mm có vân sáng của các bức xạ với bước sóng
\choice
{$0,48$ $\mu m$ và $0,56\ \mu m$}
{\True $0,40$ $\mu m$ $\text{và}$ $0,60\ \mu m$}
{$0,40\ \mu m\ \text{và}\ 0,64\ \mu m$}
{$0,45\ \mu m\ \text{và}\ 0,60\ \mu m$}

\loigiai{ 
\begin{itemize}
\item Ta có: 
$x_M=k\dfrac{\lambda D}{a}\Rightarrow \lambda =\dfrac{ax_M}{kD}=\dfrac{1,2}{k}\ \mu m\xrightarrow{{0,38\leq \lambda =\tfrac{1,2}{k}\leq 0,76}}1,58\leq k\leq 3,16\Rightarrow k=2;3
\Rightarrow \lambda =0,6\ \mu m;\ 0,4\ \mu m$
\item 
\item 
\end{itemize}
}
\end{vidu} \haidong{20}
\begin{vidu}%[3P5-6.3G][Loai3]
\nguon{QG - 2015} Trong một thí nghiệm Y-âng về giao thoa ánh sáng, khoảng cách giữa hai khe là 0,5 mm, khoảng cách từ mặt phẳng chứa hai khe đến màn quan sát là 2 m. Nguồn sáng phát ánh sáng trắng có bước sóng trong khoảng từ 380 nm đến 760 nm. M là một điểm trên màn, cách vân sáng trung tâm 2 cm. Trong các bước sóng của các bức xạ cho vân sáng tại M, bước sóng dài nhất là
\choice
{417 nm}
{570 nm}
{\True 714 nm}
{760 nm}
\loigiai{
\begin{itemize}
\item Tọa độ vân sáng: $x_{s}=k\dfrac{\lambda D}{a}=x_{M}\Rightarrow \lambda =\dfrac{a.{x}_{M}}{kD}=\dfrac{0,5.20}{k.2}=\dfrac{5}{k}\ \mu m$.
\item Theo đề bài, bước sóng nằm trong khoảng từ 380 nm đến 760 nm nên
\begin{align*}
0,38\le \lambda \le 0,76\Leftrightarrow 0,38\le \dfrac{5}{k}\le 0,76 \Rightarrow 6,6\le k\le 13,2 \Rightarrow k = 7, 8, 9, 10, 11, 12, 13.
\end{align*}
\item Bước sóng dài nhất, ứng với $k_{min}$ nên chọn $k = 7$: $\lambda _{max}=\dfrac{5}{7}=0,714\ \mu m=714\ nm.$
\end{itemize}
}
\end{vidu} \haidong{20}
\loai{Số vân tối tại M}
\begin{vidu}%[3P5-6.3K][Loai4]
Thực hiện giao thoa bằng khe Y-âng, khoảng cách giữa hai khe là 1\ mm, màn quan sát đặt song song  với mặt phẳng chứa hai khe và cách hai khe 2\ m. Chiếu sáng hai khe bằng ánh sáng trắng có bước sóng \bth{0,4\microm\leq\lambda \leq 0,75\microm}. Số bức xạ cho vân tối tại điểm M cách vân trung tâm 12\ mm là
\choice
{\True 7 bức xạ}
{5 bức xạ}
{8 bức xạ}
{6 bức xạ}
\loigiai{
\begin{itemize}
\item Tại M cách vân trung tâm 12\ mm cho vân tối nên
\begin{align*}
x_M=\pbrace{k+\dfrac{1}{2}}\dfrac{\lambda.D}{a}=12\ mm\Rightarrow \lambda=\dfrac{a.x_M}{\pbrace{k+\dfrac{1}{2}}.D}=\dfrac{1.12}{(k+0,5)2}=\dfrac{6}{k+0,5}.
\end{align*}
\item Mặt khác: \bth{0,4\microm\leq\lambda\leq 0,75\microm} nên
\begin{align*}
0,4\leq \dfrac{6}{k+0,5}\leq 0,75\Rightarrow \dfrac{6}{0,75}-0,5\leq k\leq\dfrac{6}{0,4}-0,5\Rightarrow 7,5\leq k\leq 14,5\Rightarrow k=8;\ldots;14
\end{align*}
\item Có 7 giá trị nguyên của k. 
\end{itemize}
}
\end{vidu} \haidong{20}

\begin{vidu}%[3P5-6.3G][Loai4]
Thực hiện giao thoa ánh sáng qua khe Y-âng, biết khoảng cách hai khe hẹp 1 mm, khoảng cách hai
khe tới màn là 1 m. Nguồn phát ánh sáng gồm các bức xạ đơn sắc có bước sóng trong khoảng từ $0,4\ \mu m $ đến $0,76\ \mu m $. Trên màn, tại điểm cách vân trung tâm 2 mm có vân tối của bức xạ?
\choice
{ $0,57\ \mu m $ và $0,6\ \mu m $}
{ $0,4\ \mu m $ và $0,44\ \mu m $}
{ $0,6\ \mu m $ và $0,76\ \mu m $}
{\True  $0,44\ \mu m $ và $0,57\ \mu m $}
\loigiai{
\begin{itemize}
\item Ta có: ${x_M}=k\dfrac{\lambda D}{a}\Rightarrow \lambda =\dfrac{a.x_M}{kD}=\dfrac{2,2}{k} $ (với k bán nguyên).
\item Mà $0,4<\lambda <0,76\Rightarrow 2,9<k<5,5\Rightarrow k\in \left\{3,5;4,5 \right\} $ $\Rightarrow \lambda \in \left\{0,57;0,44 \right\} $.
\item 
\item
\item   
\end{itemize}
}
\end{vidu} \haidong{20}

\loai{Bức xạ không cho vân giao thoa}
\begin{vidu}%[3P5-6.3G][Loai5]
Trong thí nghiệm giao thoa Y-âng khoảng cách hai khe là 1\ mm, khoảng cách giữa mặt phẳng chứa hai khe và màn ảnh là 1 m. Nguồn sáng S phát ánh sáng trắng có bước sóng nằm trong khoảng từ $0,38\ \mu m$ đến $0,76\ \mu m$. Tại điểm M cách vân sáng trung tâm 4\ mm bức xạ ứng với bước sóng \textbf{không} cho vân sáng là

\choice
{$\dfrac{2}{3}\ \mu m$}
{ $\dfrac{4}{9}\ \mu m$}
{ $0,5\ \mu m$}
{ \True $\dfrac{5}{7}\ \mu m$}

\loigiai{ 
\begin{itemize}
\item $x_M=k\dfrac{\lambda D}{a}\Rightarrow \lambda =\dfrac{ax_M}{kD}=\dfrac{4}{k}\ \mu m\xrightarrow{{0,38\leq \lambda =\tfrac{4}{k}\leq 0,76}}5,26\leq k\leq 10,5\Rightarrow k=6;7;8;9;10$\\
\begin{bang}[tabularx={X|X|X|X|X}]{}
$ k=6 $ & $ k=7 $ & $ k=8 $ & $ k=9 $ & $ k=10 $ \\ 
\hline 
$ \lambda =\dfrac{2}{3}\ \mu m $& $ \lambda =\dfrac{4}{7}\ \mu m $& $ \lambda =0,5\ \mu m $ & $ \lambda =\dfrac{4}{9}\ \mu m $ &$ \lambda =0,4\ \mu m $ 
\end{bang}
\item 
\item 
\item 
\item 
\end{itemize}
}
\end{vidu} \haidong{20}


